\documentclass{article}
\usepackage[margin=1in]{geometry} % Sets 1 inch margins on all sides
\usepackage{graphicx} % Required for inserting images
\usepackage[american]{babel}
\usepackage{csquotes}
\usepackage{amsmath,amssymb}
\usepackage{booktabs}
\usepackage{array}
\usepackage[backend=biber, style=apa]{biblatex}
\addbibresource{references.bib}
\usepackage[hidelinks]{hyperref}

\newcommand{\dhat}{\hat{d}}
\newcommand{\deltahat}{\hat{\delta}}
\newcolumntype{L}[1]{>{\raggedright\arraybackslash}p{#1}}

\title{Learning the Backward Pass for Test-Time Adaptation\\[4pt]
\large Research Proposal and Literature Review}
\author{Sachin Konan}
\date{October 2026}

\begin{document}

\maketitle

\noindent\emph{Status: proposal. No empirical results are claimed. Cost figures in
Section~\ref{sec:cost} are our own operation-count estimates, not measurements.}

% =====================================================================
\section{Introduction}
% =====================================================================

Meta-learning had a clear arc. In 2017, MAML made learning-to-learn a matter of
differentiating through gradient descent \parencite{finn2017maml}. By 2020, the
few-shot benchmarks it was built for were matched by a pretrained embedding and
a linear classifier \parencite{chen2019closer,tian2020rethinking,raghu2020anil},
and few-shot learning itself was absorbed by in-context learning in large
language models \parencite{brown2020gpt3}. The bi-level machinery faded.

It has now returned in a different place. Test-time training updates a model's
weights on the test stream itself, and TTT-E2E \parencite{tandon2025ttte2e}
shows that a language model whose initialization is meta-learned for test-time
gradient steps scales with context length like full attention, at constant cost
per token, with 3B-parameter models. The inner loop is no longer a few steps on
a toy task. It is the model's memory.

This puts weight on one component: the backward pass. Every gradient-based
test-time method runs one at inference. Tent \parencite{wang2021tent}
backpropagates through an entire ResNet-50 to update its normalization
parameters, under 1\% of the model. TTT-E2E restricts its inner loop to the
last quarter of its blocks, trading stored context against the cost of
backpropagating further, and its outer loop differentiates through that
backward pass, which makes training 3.4$\times$ slower than standard
pre-training at 8K context \parencite{nvidia2026ttt}. The backward pass decides how much of a network can
adapt and how expensive it is to train the adaptation. We ask:

\begin{center}
\emph{How much of the backward pass does test-time adaptation need?}
\end{center}

We propose to keep the part that is cheap and exact, and to learn the rest. The
error at the output is available in closed form. From there, a small recurrent
network runs down the depth of the model, reads the activations saved in the
forward pass, and predicts the error signal at each adapted layer for every
token. Weight updates are formed from these signals exactly as in
backpropagation, as sums of outer products $\deltahat x^\top$. The predictor is
first trained to imitate true gradients and then trained through the loss on
later data. None of these ingredients is new. Synthetic gradients, direct
feedback, and learned update rules have each been studied separately
(Section~\ref{sec:related}). Concurrent work even predicts the same per-token
signals, but it does so to parallelize training and keeps the exact backward
pass at test time \parencite{liu2026gradlev}. We study the opposite use:
predicted signals at deployment, through the full depth of a pretrained
checkpoint adapting online to a stream.

A simple cost model (Section~\ref{sec:cost}) says what to expect. The predictor
removes only the error-propagation part of the backward pass. The ceiling on
the speedup is therefore about 1.2$\times$ when a few top layers adapt,
1.5--1.75$\times$ when all layers adapt, and about 2$\times$ when only
normalization parameters adapt, as in Tent. This is why ImageNet-C is our main
benchmark, why our headline measurement is a quality--cost frontier and not a
single speedup number, and why we treat the first-order outer loop as a claim
of its own. We hope the experiments clarify which parts of backpropagation
adaptation actually depends on, whichever way they come out.

% =====================================================================
\section{Background and Related Work}
\label{sec:related}
% =====================================================================

We organize prior work around three questions: where the hidden-layer error
signal comes from, whether a backward pass runs at test time, and whether the
update rule is itself trained. Table~\ref{tab:positioning} summarizes the
answers.

% ---------------------------------------------------------------------
\subsection{Meta-Learning: From Few-Shot Benchmarks to Test-Time Training}

\subsubsection{MAML}

MAML \parencite{finn2017maml} learns an initialization from which one or a few
gradient steps on a new task's small training set give low loss on that task's
held-out examples. The outer gradient is taken through the inner gradient
steps, so it contains second derivatives of the inner loss. The paper evaluates
on few-shot classification (Omniglot, miniImageNet), sinusoid regression, and
simulated locomotion, and already reports that a first-order approximation
performs almost as well on miniImageNet. Our outer objective has the same form,
a loss measured after adaptation, but the persistent quantity we optimize is
the feedback rule and not the initialization.

\subsubsection{First-Order Variants and Feature Reuse}

Reptile \parencite{nichol2018reptile} removes second derivatives entirely: it
runs SGD on a sampled task and moves the initialization toward the resulting
weights. ANIL \parencite{raghu2020anil} then asks what MAML's inner loop is
doing. It finds that the inner loop barely changes the representations in the
network body, and that restricting adaptation to the head matches MAML. In
parallel, \textcite{chen2019closer} and \textcite{tian2020rethinking} show that
standard pre-training followed by a linear classifier matches or exceeds
meta-learners on the same benchmarks. Together these results explain why the
field moved on: on narrow task distributions a good representation was
sufficient, and the inner loop was not doing the work. We take this as a design
rule. Every claim in this proposal is paired with a simple baseline that could
make the learned component redundant (Section~\ref{sec:baselines}).

\subsubsection{In-Context Learning}

GPT-3 \parencite{brown2020gpt3} showed few-shot learning emerging from
next-token prediction at scale, with no weight updates.
\textcite{vonoswald2023icl} connect this back to optimization: a linear
self-attention layer can implement a step of gradient descent on a regression
loss in its forward pass. On this view in-context learning and explicit inner
loops belong to one family. The choice between them is about where the adapted
state lives, in activations or in weights, and what it costs to update.

% ---------------------------------------------------------------------
\subsection{Synthetic Gradients}

\subsubsection{Decoupled Neural Interfaces}

DNI \parencite{jaderberg2017dni} attaches a small model to a layer's output $h$
that predicts $\partial L/\partial h$, so that the layers below can update
without waiting for the rest of the network. The predictor is trained by
regression onto the gradient that later arrives from above, which may itself be
a synthetic gradient from the next module. On feed-forward networks (MNIST,
CIFAR-10) training succeeds with some loss of accuracy, part of which is
recovered when the predictor is conditioned on the label. In recurrent
networks, synthetic gradients extend the effective horizon of truncated
backpropagation through time. DNI is the direct ancestor of our imitation
phase. Two differences matter. DNI aims to remove update locking, not to reduce
computation: each module still backpropagates locally. And its predictor chases
a network trained from scratch, whereas ours is fit around a fixed pretrained
checkpoint.

\subsubsection{Analysis of Synthetic Gradients}

\textcite{czarnecki2017understanding} study what synthetic gradients do to
optimization. They prove convergence for linear and deep linear models, show
empirically that the learned layer-wise representations differ from those of
backpropagation, and place feedback alignment and direct feedback alignment in
the same framework as special cases with fixed predictors. These results
motivate our correctness checks on linear models. They do not cover a
nonlinear, recurrent predictor trained through an outer loss.

\subsubsection{BP($\lambda$)}

\textcite{pemberton2024bplambda} train synthetic-gradient models fully online
in recurrent networks, using eligibility traces in the manner of TD($\lambda$)
so that producing targets needs no backpropagation through time. The work makes
the bias of bootstrapped targets explicit. It concerns credit assignment over
time; ours is over depth.

\subsubsection{Adaptive Gradient Prediction}

ADA-GP \parencite{janfaza2023adagp} alternates between phases that run true
backpropagation while training a gradient predictor and phases that skip
backpropagation and use predicted gradients, with phase lengths chosen
adaptively. The authors note that predicting all the time would remove the
backward pass entirely but degrades accuracy. This is prior art for a hybrid
variant of our method with periodic exact refreshes, and it suggests the
matched control: exact updates at the same reduced frequency.

\subsubsection{When Synthetic Gradients Beat Backpropagation}

\textcite{zhang2026backprop} view backpropagation as a Monte Carlo estimator of
the expected gradient and synthetic gradients as a bootstrapped one, and
characterize when the latter has lower mean squared error: some source of
randomness in the feedback, or structural sparsity, is required. Experiments
are on contextual bandits and reinforcement learning. The paper also notes that
when exact backpropagation is available, prior work finds synthetic gradients
do not match it without additional mechanisms. The implication for us is
narrow. A predicted signal can beat the minibatch gradient only as an estimate
of the \emph{expected} gradient, which matters on noisy streams and not for a
fixed observed batch.

\subsubsection{A Recent Negative Result}

A 2026 preprint (arXiv:2608.21386) reports in an appendix that per-layer
gradient predictors in a Transformer became redundant once local auxiliary
losses were present, and that the cosine similarity between predicted and true
gradients fell from about 0.99 at initialization to 0.2--0.5 near convergence.
We have read only an excerpt of this work. If the decay also holds around a
converged checkpoint under distribution shift, our first claim fails, so we
measure it before building anything else (Section~\ref{sec:gates}).

% ---------------------------------------------------------------------
\subsection{Feedback Without Weight Transport}

\subsubsection{Feedback Alignment}

\textcite{lillicrap2016random} replace the transposed forward weights in the
backward pass with fixed random matrices. Learning still works, because the
forward weights come to align with the feedback weights. The result established
that useful error signals do not require the exact transpose.

\subsubsection{Direct Feedback Alignment}

DFA \parencite{nokland2016dfa} goes further and sends the output error to every
hidden layer through its own fixed random projection, with no layer-to-layer
backward pass at all. \textcite{launay2020dfa} show that DFA trains modern
architectures in neural view synthesis, recommender systems, graph learning,
and language modeling with performance close to backpropagation; convolutional
networks remain a known difficult case. DFA is the essential cheap competitor.
It has access to the same output error as our method and costs one projection
per layer. Our predictor must beat it in both a fixed-random and a
regression-fit form.

\subsubsection{Learned Feedback Weights}

\textcite{akrout2019weighttransport} learn the feedback weights instead of
fixing them, with rules that drive them toward the transposes of the forward
weights, and reach accuracy close to backpropagation on ImageNet with ResNet-18
and ResNet-50. Learned feedback can therefore work at ImageNet scale, although
here the backward pass costs as much as backpropagation.

\subsubsection{Local Losses}

A different way to avoid deep backward passes is to give each layer or block
its own loss. \textcite{nokland2019local} and \textcite{belilovsky2020decoupled}
train deep convolutional networks with local auxiliary objectives and approach
end-to-end accuracy. Local losses are the main alternative to predicted
gradients, and the test-time-training literature has its own version
(Section~\ref{sec:tttseq}).

% ---------------------------------------------------------------------
\subsection{Learned Update Rules}

\subsubsection{Learned Optimizers}

\textcite{andrychowicz2016learning} train an LSTM that maps each coordinate's
gradient history to an update, by backpropagating through unrolled
optimization. To keep this tractable they ignore the dependence of the
optimizee's gradients on the optimizer's parameters, which drops the
second-derivative terms. The recurrence over optimization steps is the
precedent for memory across batches. The optimizer consumes true gradients, so
it says nothing about skipping the backward pass.

\subsubsection{Meta-Learned Update Rules}

\textcite{metz2019metalearning} meta-learn a local update rule for unsupervised
representation learning. The rule includes its own learned backward weights,
untied from the forward weights, and the meta-objective is few-shot
classification from the learned features. The rule transfers across
architectures and datasets. It is an early instance of learning the feedback
pathway through an outer objective.

\subsubsection{Variable Shared Meta Learning}

\textcite{kirsch2021vsml} replace each weight with a tiny recurrent network
whose parameters are shared across the whole model, with messages passing
forward and backward. Such a system can implement backpropagation in its
forward dynamics, and it can be meta-trained into learning algorithms that
differ from backpropagation. It is the closest conceptual precedent for a
learned backward algorithm. It does not address cost.

\subsubsection{Feedback and Local Plasticity}

\textcite{lindsey2020flp} meta-learn feedback weights, the forward
initialization, and plasticity rates so that a network learns online with local
Hebbian updates and no gradient computation in the inner loop. The meta-trained
networks match or exceed a gradient-based online meta-learning baseline on
regression and classification, with the largest gains in continual learning.
This is the closest prior work to our outer-loop phase: feedback trained
through the performance of the adapted network. The differences are scale and
setting. Their networks are small and start from a meta-learned initialization
on few-shot tasks. Ours are pretrained ImageNet-scale models adapting to a
stream, with the exact shortcut path retained and gradient imitation as the
starting point.

\subsubsection{MEND}

MEND \parencite{mitchell2022mend} edits large language models by transforming
fine-tuning gradients with small learned networks. It relies on the same
identity as our method: for a linear layer the gradient on one example is the
rank-one product $\delta x^\top$, so the editor networks act on $\delta$ and
$x$ and never on the full weight matrix. MEND takes the true $\delta$ from a
backward pass and learns to transform it. We predict $\delta$.

\subsubsection{Learned Gradient Generation for Adaptation}

MGTTA \parencite{deng2025mgtta} learns an optimizer for test-time adaptation
that stores historical gradient information in a memory layer and refines the
update. It is the nearest history-based method in our target setting. It
consumes gradients computed by backpropagation, so it is a quality baseline and
not a cost baseline. It is also the template for a control this proposal needs:
a learned rule fed with exact signals.

% ---------------------------------------------------------------------
\subsection{Test-Time Adaptation in Vision}

\subsubsection{Test-Time Training}

\textcite{sun2020ttt} train a network with a shared encoder, a classification
head, and a self-supervised head that predicts image rotation. At test time the
encoder is updated on the self-supervised loss of the incoming sample, either
independently per sample or online over a stream. This improves robustness on
CIFAR-10-C and ImageNet-C and defines the label-free setting.

\subsubsection{Normalization Statistics}

\textcite{schneider2020covariate} show that re-estimating batch normalization
statistics on test data removes a large part of the corruption error on
ImageNet-C with no gradient step. It is the forward-only baseline that any
backward-free method must beat.

\subsubsection{Tent}

Tent \parencite{wang2021tent} needs no change to training. It minimizes the
entropy of the model's predictions on each test batch and updates only the
affine parameters of the normalization layers, with normalization statistics
taken from the batch. Few parameters change, but the gradient with respect to
an early normalization layer requires a backward pass through nearly the whole
network. Tent is our primary exact baseline, and its structure is why the cost
ceiling of our method is highest here (Section~\ref{sec:cost}).

\subsubsection{EATA}

EATA \parencite{niu2022eata} adds two things to entropy minimization. It adapts
only on samples that are reliable and non-redundant, which reduces the number
of backward passes, and it regularizes important weights to limit forgetting.
It is a cost baseline of a different kind: fewer exact backward passes, not
cheaper ones.

\subsubsection{CoTTA}

CoTTA \parencite{wang2022cotta} defines continual test-time adaptation, where
the target distribution changes over a long stream with no reset. It uses a
weight-averaged teacher, augmentation-averaged pseudo-labels, and stochastic
restoration of weights to the source model. Its protocol on CIFAR-10-C,
CIFAR-100-C, and ImageNet-C is our long-horizon stability test.

\subsubsection{SAR}

\textcite{niu2023sar} study adaptation under mixed shifts, small batches, and
label shift. Batch-normalized models are unstable in these conditions; group-
and layer-normalized models are more stable but can still collapse. The remedy
is to filter high-entropy samples and to minimize entropy with a
sharpness-aware step. This matters for us because a rule trained on future loss
could learn such caution, or could inherit the collapse.

\subsubsection{Forward-Only Adaptation}

FOA \parencite{niu2024foa} adapts a vision Transformer without any backward
pass. It learns an input prompt by derivative-free optimization, with a fitness
that combines prediction entropy and activation statistics, and it shifts the
final-layer activations. It runs on quantized models, where backpropagation is
unavailable. FOA shows there is demand for backward-free adaptation and is the
closest competitor on that axis. It pays with several forward passes per batch.

% ---------------------------------------------------------------------
\subsection{Test-Time Training for Sequences}
\label{sec:tttseq}

\subsubsection{Dynamic Evaluation}

\textcite{krause2018dynamic} adapt a language model's weights to the recent
history of the test sequence by gradient descent during evaluation, with a
decay back toward the trained weights, and improve perplexity on standard
benchmarks; \textcite{krause2019dynamic} extend this to Transformers. The model
is not trained for the adaptation it will undergo. Dynamic evaluation is the
language analog of Tent and our bridge benchmark.

\subsubsection{TTT Layers}

\textcite{sun2024tttlayers} make the hidden state of a sequence layer the
weights of a small inner model. The inner model is updated per token by a
gradient step on a self-supervised reconstruction loss whose inputs and targets
are learned in the outer loop. Because the loss is local to the layer, the
inner gradient needs no backward pass through depth. Later work scales the idea
with large-chunk updates \parencite{zhang2025lact} and adds momentum and
forgetting to the inner update \parencite{behrouz2025titans}.

\subsubsection{TTT-E2E}

TTT-E2E \parencite{tandon2025ttte2e} replaces the layer-wise loss with the
next-token loss at the end of the network, and meta-learns the initialization
so that the loss \emph{after} test-time updates is low. The architecture is a
standard Transformer with sliding-window attention over 8K tokens. The inner
loop takes one step per 1K tokens and updates only the MLPs in the last quarter
of the blocks. With 3B-parameter models it scales with context length in the
same way as full attention, with inference 2.7$\times$ faster than full
attention at 128K context.

Three of its findings set up this proposal. First, the end-to-end signal is
better than the local one: at 760M parameters and 8K context, the loss falls
from 2.819 to 2.806 when the layer-wise loss is replaced by the next-token
loss. Second,
the cost of preparing upstream gradients shapes the design: the authors update
fewer, larger layers because backpropagating to many small ones is expensive.
Third, the outer loop needs gradients of gradients, and the reported training
cost is 3.4$\times$ that of standard pre-training at 8K context, attributed to
missing second-order support in the attention kernel
\parencite{nvidia2026ttt}. Our method targets the space between the first
finding and the other two: an end-to-end error signal at close to local cost,
with a first-order outer loop.

\subsubsection{GradLev: Costate Prediction}
\label{sec:gradlev}

GradLev \parencite{liu2026gradlev} is the closest work to this proposal. Its
problem is training, not deployment. Token-level test-time training is serial:
the weights used at token $t$ depend on the gradient at token $t-1$. GradLev
observes that once the \emph{costates}, the gradients at the outputs of the
adapted layers, are given, online gradient descent is an affine recurrence in
the weights. Parallel scans can then evaluate both the forward pass and the
transposed backward pass. An auxiliary causal network predicts the costates
for all tokens at once, from features below the adapted layers and the
embedding of the observed target. A real backward pass through the resulting
trajectory supplies exact costates, and a consistency loss pulls the
predictions toward them. At exact consistency the parallel trajectory equals
the serial one. At deployment the predictor is discarded, and the model updates
from exact per-token gradients. The experiments are preliminary. A 300M
language model that adapts only its final MLP improves validation loss by about
0.002 nats over chunked updates and 0.003 over a static model, and the
differences in downstream accuracy lie within the reported intervals.

Three things carry over to our setting. The mechanism overlaps with ours:
per-token error signals predicted by an auxiliary network, written as outer
products, with the predictor trained on the trajectory that its own predictions
produce (our Phase~B). GradLev's training is first-order as well. It stops
gradients through the writes, which recovers a first-order MAML gradient at
exact consistency. And predicted costates remove the serial dependence between
tokens, which our exact seed does not: the output error at token $t$ depends on
the adapted output.

The differences define our contribution. GradLev keeps the backward pass at
test time; we remove it. GradLev adapts one MLP and names a predictor that
stays accurate across layers as an open problem; the recurrence over depth
(Section~\ref{sec:recurrence}) is our answer to it. GradLev's predictor must
match the gradient, because the deployed model uses the gradient; ours may
depart from it when that lowers future loss. Finally, by our estimate, with one
adapted MLP their two-block predictor costs more than the single-layer backward
pass it would replace, so there is no reason to keep it at deployment. That
changes once adaptation is deep, which is the regime we target.

% ---------------------------------------------------------------------
\subsection{Positioning}

\begin{table}[ht]
\centering
\small
\begin{tabular}{L{3.0cm} L{4.4cm} L{3.6cm} L{3.4cm}}
\toprule
Method & Hidden-layer signal & Backward pass at test time & Learned rule \\
\midrule
Tent & exact, entropy loss & full network & no \\
EATA & exact, entropy loss & full, on selected samples & no \\
FOA & none (derivative-free) & none & no \\
DNI & predicted, bootstrapped & within each module & predictor only \\
DFA & fixed random projection of output error & none & no \\
FLP & meta-learned feedback & none & yes \\
MGTTA & exact, then transformed & full network & yes \\
TTT layers & exact, local loss & one layer & outer loop \\
TTT-E2E & exact, end loss & last quarter of blocks & yes, second-order \\
GradLev & predicted in training, exact at test time & adapted layers, every token & yes, first-order \\
\midrule
This proposal & predicted from the exact output error & none beyond the output seed & yes, first-order \\
\bottomrule
\end{tabular}
\caption{Where each method gets its hidden-layer error signal and what it pays
for it. The proposal occupies the empty cell: an end-loss signal with no
backward pass through depth, trained through its effect on later data.}
\label{tab:positioning}
\end{table}

% =====================================================================
\section{Proposed Method}
\label{sec:method}
% =====================================================================

\subsection{Setting}

A pretrained network processes a stream of batches $C_1, C_2, \dots$ A chosen
subset $W$ of its parameters, the fast weights, is updated online; everything
else stays at the checkpoint. Each batch is scored with the current weights
\emph{before} it is used for an update. The predictor's parameters $\phi$ are
fixed at deployment. We train one predictor per checkpoint and architecture;
transfer across checkpoints is a separate question.

\subsection{Local Gradient Identities}

Let $n$ index examples and positions (tokens or spatial locations). For a
linear map $z = Wx$ with error signal $\delta = \partial L/\partial z$,
\begin{equation}
\nabla_W L = \sum_n \delta_n x_n^\top .
\label{eq:linear}
\end{equation}
For a normalization layer $y = \gamma \odot \tilde{x} + \beta$ with
$\delta = \partial L/\partial y$ and normalized input $\tilde{x}$,
\begin{equation}
\nabla_\gamma L = \sum_n \delta_n \odot \tilde{x}_n, \qquad
\nabla_\beta L = \sum_n \delta_n .
\label{eq:norm}
\end{equation}
Convolutions follow \eqref{eq:linear} with input patches in place of $x$. In
both cases the forward quantity is already in memory, and the only missing
quantity is $\delta$, the costate in the terminology of GradLev
\parencite{liu2026gradlev}. The method predicts $\delta$ for every token and leaves
\eqref{eq:linear}--\eqref{eq:norm} untouched. It does not predict weight
matrices or low-rank factors. Per-token signals keep the procedure identical to
backpropagation with one term replaced, which is what makes the comparison with
exact adaptation clean.

\subsection{The Exact Seed}

The error at the logits $z$ is available in closed form for each setting we
study. With softmax output $p$, $n$ scored predictions, and entropy
$H = -\sum_k p_k \log p_k$:
\begin{align}
\text{labels revealed (cross-entropy):} \quad &
  \partial L/\partial z = (p - y)/n, \\
\text{label-free (entropy, as in Tent):} \quad &
  \partial H/\partial z_j = -\,p_j\,(\log p_j + H), \\
\text{language modeling:} \quad &
  \text{cross-entropy against the observed next token.}
\end{align}
The seed is propagated exactly through the frozen head and any final pooling or
normalization, giving the error $d_{L+1}$ at the top of the residual stream.
This is the only true gradient the deployed method uses.

\subsection{A Recurrence over Depth}
\label{sec:recurrence}

Write block $l$ as $h_{l+1} = \sigma_l\!\left(P_l h_l + F_l(h_l)\right)$, where
$P_l$ is the shortcut (identity or projection), $F_l$ the residual branch, and
$\sigma_l$ the identity for pre-activation and pre-norm designs or a ReLU for
post-activation ResNets. With $d_l = \partial L/\partial h_l$ and $m_l$ the
derivative of $\sigma_l$ at its input, backpropagation gives
\begin{equation}
d_l = P_l^\top (m_l \odot d_{l+1}) + J_{F_l}^\top (m_l \odot d_{l+1}).
\label{eq:exact}
\end{equation}
The first term is cheap and exact. We keep it and replace the second with a
low-rank learned correction driven by a small state $s_l \in \mathbb{R}^r$ per
position:
\begin{align}
s_{L+1} &= E^\top d_{L+1}, \qquad \dhat_{L+1} = d_{L+1}, \\
s_l &= s_{l+1} + B\left(\varphi\!\left(C_l^\top a_l + e_l\right) \odot A\, s_{l+1}\right),
\label{eq:cell} \\
\dhat_l &= P_l^\top (m_l \odot \dhat_{l+1}) + U_l\, s_l .
\label{eq:pred}
\end{align}
Here $a_l$ is a forward activation of block $l$ that is already saved (the
input to the adapted layer), $e_l$ is a learned layer embedding, $\varphi$ is a
pointwise nonlinearity, $A$ and $B$ are $r \times r$ matrices shared across
depth, and $C_l$, $U_l$, $E$ project between the model width and $r$. The cell
is multiplicative by design: the true term in \eqref{eq:exact} is linear in the
incoming error and gated by the forward activations, and \eqref{eq:cell} has
the same form. $U_l$ is initialized to zero, so the rule starts as
\emph{shortcut-only} feedback, a transparent baseline.

\paragraph{Taps.} An adapted parameter needs the error at its own output. For a
layer that writes directly into the residual sum (the last normalization layer
of a ResNet bottleneck, or the output projection of a Transformer MLP), that
error is $m_l \odot \dhat_{l+1}$ and comes for free from \eqref{eq:pred}. For a
parameter inside a branch (other normalization layers; LayerNorm at the input
of a Transformer sublayer) we add a linear head $\deltahat = V s_l$. The first
experiments use only the first kind.

\paragraph{Mixing across positions.} Equations
\eqref{eq:cell}--\eqref{eq:pred} treat positions independently, which makes
them a batch dimension. The true error is not independent across positions: in
a causal Transformer, position $t$ receives error from the losses at later
positions that attend to it, and in a ConvNet from neighboring locations. We
add a cheap mixing step on $s$, a depthwise $3\times 3$ convolution for
ConvNets and an anti-causal depthwise 1-D convolution for sequences. The
position-independent predictor is an ablation.

\paragraph{Parallelism.} Positions are processed in parallel. The recurrence is
sequential over $L$ blocks, each step a pair of small matrix products. With an
exact seed, updates are applied once per chunk (a batch of images, or a
mini-batch of tokens as in TTT-E2E): the output error at a token depends on the
adapted output, so updating after every token would make the forward pass
serial. GradLev removes this dependence by predicting costates from features
that do not depend on the fast weights, after which token-level updates are
parallel scans \parencite{liu2026gradlev}. This gives two variants: an exact
seed with chunked updates, our default, and a predicted seed with token-level
updates. We compare them on the language task. The sums in \eqref{eq:linear}--\eqref{eq:norm} are batched matrix
products. In JAX: \texttt{einsum} for the updates, \texttt{lax.scan} over
depth, and \texttt{vmap} over sequences when each sequence carries its own fast
weights.

\subsection{Training the Predictor}

\paragraph{Phase A: imitation.} Run exact adaptation episodes on training
streams and record the true error at every tap before each update. Train
$\phi$ to regress these targets with a per-layer normalized squared error,
feeding the exact incoming error at first. Report direction and magnitude
separately; cosine alone does not calibrate a step.

\paragraph{Phase B: self-generated states.} Roll the recurrence out from the
seed with its own predictions as inputs, and let the predicted updates move the
working weights. Recompute true targets on the states the predictor actually
reaches. This corrects two mismatches: error compounding down the depth, and
drift of the weights over a stream.

\paragraph{Phase C: future loss.} Unroll $K$ updates and train $\phi$ on the
loss of later batches:
\begin{equation}
J(\phi) = \mathbb{E}\Big[\sum_{t=1}^{K} \ell\big(C_t; W_t\big)\Big],
\qquad W_{t+1} = W_t - \eta\, g_\phi(W_t, C_t),
\end{equation}
where $g_\phi$ denotes the predicted update and each $C_t$ is scored before
$W_t$ is updated on it. Two properties are worth stating. First, $\ell$ can be
the supervised loss even when the seed is label-free: labels are available on
the predictor's training streams, so an entropy-driven rule can be trained to
reduce true error on later data. Second, the outer gradient is first-order.
$g_\phi$ is a composition of a forward pass, a closed-form seed, the recurrence
\eqref{eq:cell}--\eqref{eq:pred}, and the sums
\eqref{eq:linear}--\eqref{eq:norm}. It contains no call to reverse-mode
differentiation, so $\nabla_\phi J$ is one ordinary reverse pass through the
unroll. With an exact inner gradient the same objective requires
reverse-over-reverse. First-order training alone is not new: GradLev obtains it
by stopping gradients through the writes, which approximates the
meta-gradient. Here the writes are an explicit function of $\phi$, so the
first-order pass is the exact gradient of the learned rule. We start at $K=2$, keep an imitation term with a decaying
weight, and always compare three variants: imitation only, outer training from
random initialization, and imitation followed by outer training.

\subsection{Cost Model}
\label{sec:cost}

Let $F$ be the cost of the forward pass, $V$ the cost of propagating the error
through the adapted depth (the vector--Jacobian products), and $G$ the cost of
forming parameter gradients from $\delta$ via
\eqref{eq:linear}--\eqref{eq:norm}. Exact adaptation costs $F + V + G$. The
learned rule costs $F + G + \Pi$, where $\Pi$ is the predictor. It removes $V$
and nothing else. For a linear or convolutional layer, each of the three costs
is about one forward-equivalent of that layer, because the sum over tokens in
\eqref{eq:linear} is a matrix product of the same size as the forward one. All
terms scale linearly with batch size and sequence length, so the ratio does not
improve with scale.

\begin{table}[ht]
\centering
\small
\begin{tabular}{L{7.6cm} c c c c}
\toprule
Configuration & $F$ & $V$ & $G$ & Ceiling $\frac{F+V+G}{F+G}$ \\
\midrule
ResNet, 9 two-conv blocks; final conv of top 3 blocks & 18.2 & 4 & 3 & 1.19$\times$ \\
Same network; final conv of all 9 blocks & 18.2 & 16 & 9 & 1.59$\times$ \\
Transformer; MLP output projection, last quarter of blocks & 3 & 0.75 & 0.25 & 1.23$\times$ \\
Transformer; MLP output projection, all blocks & 3 & 3 & 1 & 1.75$\times$ \\
Any network; normalization parameters only, all layers & 1 & $\approx 1$ & $\approx 0$ & $\approx 2\times$ \\
\bottomrule
\end{tabular}
\caption{Speedup ceilings with a free predictor ($\Pi = 0$). ResNet rows are in
units of one $3\times 3$ convolution; Transformer rows are in units of one MLP
matrix product per block, with attention projections counted as one unit and
attention-score computation ignored. These are operation counts, not wall-clock
measurements. The ceiling is highest when the parameter gradient is cheap.}
\label{tab:cost}
\end{table}

Three consequences follow from Table~\ref{tab:cost}. Adapting only a few top
layers leaves almost nothing to save. Adapting normalization parameters, where
$G \approx 0$, is the most favorable case, which is the Tent setting. And the
exact baseline can choose its depth too: by this count, predicted signals on
all Transformer blocks cost the same as exact signals on the last quarter. The
honest comparison is therefore a frontier, quality against cost at several
adapted depths for both methods, not a single ratio. The predictor itself is
small: at $r = d/16$ and a $4\times$ MLP, one step of
\eqref{eq:cell}--\eqref{eq:pred} costs about $0.3\,d^2$ per token, under 3\% of
the roughly $12\,d^2$ vector--Jacobian product of a Transformer block. A
full-width recurrent cell would cost more than the pass it replaces.

% =====================================================================
\section{Experimental Plan}
\label{sec:experiments}
% =====================================================================

\subsection{Benchmarks}

\begin{table}[ht]
\centering
\small
\begin{tabular}{L{3.9cm} L{3.5cm} L{2.0cm} L{5.2cm}}
\toprule
Benchmark & Models & Seed & Role \\
\midrule
CIFAR-10-C, CIFAR-100-C \parencite{hendrycks2019imagenetc}
  & standard public checkpoints & entropy & development loop; all ablations run here first \\
ImageNet-C, 15 corruptions, severity 5
  & ResNet-50, ViT-B/16, ConvNeXt \parencite{he2016resnet,dosovitskiy2021vit,liu2022convnext}
  & entropy & main result; quality--cost frontier \\
Continual ImageNet-C \parencite{wang2022cotta}
  & same & entropy & long streams with no reset; stability \\
ImageNet-R, ImageNet-Sketch, ImageNetV2
  \parencite{hendrycks2021imagenetr,wang2019sketch,recht2019imagenetv2}
  & same & entropy & shifts never seen by the predictor \\
WikiText-103 \parencite{merity2017wikitext}, enwik8
  & small Transformer & next token & dynamic evaluation; bridge to language \\
Long-context language modeling
  & small Transformer with sliding-window attention & next token
  & TTT-E2E-style comparison; outer-loop cost \\
\bottomrule
\end{tabular}
\caption{Benchmarks. Each uses a pretrained checkpoint, an online stream, and
an exact seed at the output.}
\label{tab:benchmarks}
\end{table}

Training ImageNet from scratch with predicted gradients is out of scope. The
predictor would have to track weights that move far from any checkpoint, which
is the original synthetic-gradient problem, and the outer objective cannot be
unrolled through a full training run. Reinforcement learning is also out of
scope: policy networks are small, so the cost argument does not apply, and the
relevant claim there is about gradient noise \parencite{zhang2026backprop}.

\subsection{Predictor Training Data}

The predictor is a trained component, so the data that trains it is training
data for the whole system. Two rules follow.

\emph{No test corruption in training.} ImageNet-C defines 15 test corruption
types and four additional held-out types (speckle noise, Gaussian blur,
spatter, saturate) intended for validation
\parencite{hendrycks2019imagenetc}. Predictor training streams use the four
held-out types plus generic photometric and geometric shifts, applied to
training images. None of the 15 test types appears at any severity. Stream
order, shift schedule, and all hyperparameters are fixed before the test
streams are run.

\emph{Check the backbone-seen confound.} Public ImageNet checkpoints were
trained on the full training set, so the predictor's training images have been
seen by the backbone, and their error statistics may differ from those of test
images. On CIFAR we control this directly, by retraining the backbone with a
held-out split and comparing predictors trained on seen and unseen images. If
the difference matters, the ImageNet experiments split the validation set by
image and report on the held-out half, flagged as non-standard.

\subsection{Baselines and Controls}
\label{sec:baselines}

\begin{table}[ht]
\centering
\small
\begin{tabular}{L{5.6cm} L{9.4cm}}
\toprule
Method & What it controls for \\
\midrule
No adaptation & whether adaptation helps at all \\
Normalization statistics only \parencite{schneider2020covariate} & what a single forward pass already buys \\
Tent, all normalization parameters & the standard exact method \\
Tent, restricted to our taps & exact signals at the same adapted parameters \\
EATA, SAR & exact methods with sample selection and stability mechanisms \\
FOA & backward-free adaptation by other means \\
Shortcut-only feedback ($U_l = 0$) & what the exact identity path gives without learning \\
DFA, fixed random and regression-fit & whether a recurrence beats one projection per layer \\
Imitation only (Phases A--B) & whether training on future loss adds anything \\
Outer training from random initialization & whether imitation is a useful starting point \\
\textbf{Learned rule on exact signals} & separates the signal source from the learned update rule \\
Exact adaptation at depth $\tfrac18, \tfrac14, \tfrac12, 1$ & the exact quality--cost frontier \\
Exact updates at matched reduced frequency & whether predicting between exact steps beats skipping them \\
Local-loss TTT layer (language) & the existing cheap alternative to a deep backward pass \\
GradLev as published (language) & exact per-token gradients at test time; the reference for token-level adaptation \\
GradLev with its predictor kept at deployment & predicted costates without an exact seed or a recurrence over depth \\
Predicted seed with token-level updates & whether the exact seed is worth giving up token-level parallelism \\
\bottomrule
\end{tabular}
\caption{Baselines. The row in bold is the control missing from most learned-feedback
work: Phase~C changes both the signal and the rule, and without this row a
gain cannot be attributed.}
\label{tab:baselines}
\end{table}

\subsection{Measurements and Decision Gates}
\label{sec:gates}

We measure three things separately. \emph{Signal quality}: per-layer cosine and
norm ratio against the true error, as a function of depth, on states the
predictor reaches by itself, and the same for the induced parameter gradient.
\emph{Adaptation quality}: error of predictions made before their own update,
averaged over the stream; recovery after a shift; clean accuracy after
adaptation. When exact adaptation improves on the frozen model, we also report
$R = (\text{err}_\text{frozen} - \text{err}_\text{learned}) /
(\text{err}_\text{frozen} - \text{err}_\text{exact})$. \emph{Cost}: operation
counts, wall-clock per batch, and peak memory, with the predictor and all
caches included.

Effect sizes differ sharply between the two domains. On ImageNet-C, adaptation
moves error by many points. On language, GradLev's gain over chunked updates is
about 0.002 nats \parencite{liu2026gradlev}, and the gain TTT-E2E obtains from
the end-to-end loss is 0.013 \parencite{tandon2025ttte2e}. Language
comparisons therefore use three or more training seeds, paired evaluation on a
fixed token set, and intervals over seeds.

The gates are fixed before any test stream is run.

\begin{enumerate}
\item \textbf{Correctness.} With true errors inserted, the update routine
reproduces the exact baseline. The outer gradient matches finite differences on
a small smooth model.
\item \textbf{Predictability.} At the pretrained checkpoint under shift, a
regression-fit predictor reaches useful cosine with the true error at depth.
This is a week-one measurement. If it fails, the project stops or narrows.
\item \textbf{Adaptation.} The learned rule beats no adaptation,
normalization-statistics adaptation, shortcut-only feedback, and DFA.
\item \textbf{Attribution.} Against the learned rule on exact signals, the
predicted signal retains most of the gain. Otherwise the result is about the
rule and not the signal.
\item \textbf{Frontier.} At one or more cost levels, the learned curve lies
below the exact curve by more than seed variance. This replaces a fixed speedup
threshold.
\item \textbf{Training cost.} On the language task, wall-clock per outer step
is lower than for the second-order outer loop at equal model, context, and
adapted layers.
\end{enumerate}

\subsection{Schedule}

\begin{table}[ht]
\centering
\small
\begin{tabular}{L{2.0cm} L{13.0cm}}
\toprule
Weeks & Work \\
\midrule
1 & JAX infrastructure, residual-MLP correctness checks, cost profiler; predictability measurement on CIFAR-10-C \\
2 & Reproduce exact baselines on CIFAR-10-C and CIFAR-100-C \\
3--4 & Predictor Phases A--B; shortcut-only and DFA controls \\
5--6 & Phase C; learned rule on exact signals; first frontier on CIFAR-C \\
7--8 & ImageNet-C with ResNet-50, then ViT-B/16; continual protocol \\
9 & Natural shifts; ConvNeXt \\
10--11 & Dynamic evaluation; TTT-E2E-style small model; outer-loop cost \\
12 & Replication across seeds and checkpoints; write-up \\
\bottomrule
\end{tabular}
\caption{Indicative schedule for one researcher. Work stops or narrows at a
failed gate.}
\label{tab:schedule}
\end{table}

% =====================================================================
\section{Risks and Open Questions}
% =====================================================================

\paragraph{Gradients may be hard to predict near convergence.} The negative
result discussed in Section~\ref{sec:related} points this way. Distribution
shift may restore structure to the error signal, but this is an empirical
question, and it is the first one we answer.

\paragraph{A rank-$r$ correction may be too weak.} The true branch term in
\eqref{eq:exact} is full-rank. We sweep $r$ and report signal quality against
predictor cost. If useful quality needs $r$ close to $d$, the cost argument
fails.

\paragraph{The equal-cost comparison is demanding.} TTT-E2E found that updating
12 of 24 layers performed about the same as updating 6
\parencite{tandon2025ttte2e}. If more adapted layers do not help even with
exact gradients, noisier signals on more layers will not either, and the
frontier will show it.

\paragraph{Entropy is a weak seed.} Entropy minimization is known to collapse
on long or mixed streams \parencite{niu2023sar,wang2022cotta}. Training the
rule on supervised future loss is meant to counter this. It could instead
inherit the failure, or overfit to the shift types seen in training.

\paragraph{Cross-position error is only approximated.} The mixing step is
cheap, and the true dependence through attention is not. The
position-independent ablation measures how much this matters.

\paragraph{Memory.} When each sequence carries its own fast weights, the update
tensors are weight-sized per layer per sequence, and the outer loop holds $K$
copies. This bounds batch size and unroll length before computation does.

\paragraph{Concurrent work may close the gap.} GradLev becomes a backward-free
method if its predictor is kept at deployment
(Section~\ref{sec:gradlev}). Its authors have no reason to do this while
adapting one layer, but deeper adaptation is their stated next step. Our vision
experiments lie outside that line of work, and on language we run the
kept-predictor variant ourselves as a baseline.

\paragraph{The predictor is tied to a checkpoint.} Its offline training cost is
reported and must be amortized over deployment. Whether one predictor serves
several checkpoints of an architecture is left for later.

% =====================================================================
\section{Summary}
% =====================================================================

The proposal tests one idea: at test time, the error at the output is exact and
cheap, and the rest of the backward pass can be replaced by a small learned
recurrence over depth that emits per-token error signals. The idea is assembled
from parts that exist. What is new is where the predicted signal is used, at
deployment and through the full depth of a pretrained checkpoint adapting to a
stream, and the evaluation, which separates three claims that are
usually reported together: that the signals can be learned, that they support
adaptation, and that they cost less. The cost model bounds the third claim in
advance, at about 2$\times$ in the best case. Within that bound, the result we
are after is a better quality--cost frontier on ImageNet-C and an outer loop
that needs no gradients of gradients.

\printbibliography

\end{document}
